\section{Introduction}

\subsection{The problem that will not go away}
Why do you feel pain? This is the simplest question in the world. It is also the hardest. We all know what pain is. We have all felt it. But why do we feel it? We know that nociceptors send signals to the thalamus, that the cortex processes them and that endorphins are released. All of this is a functional description. It tells us what the brain does, but it does not tell us why you feel it.

Think of a computer. It can process information, it can recognise an image, it can play chess. But it feels nothing. It has no subjective experience. Nothing hurts it. Why is the brain different? This is the hard problem of consciousness, formulated by David Chalmers in 1995. Neuroscience has made enormous progress since then, but the hard problem has remained intact. No one has managed to explain why matter becomes experience. In 1998 Christof Koch bet Chalmers that a neural correlate of consciousness would be identified within 25 years. In 2023 he conceded that the bet was lost.

\subsection{Dark matter --- the exact numbers}
The universe has a composition we know fairly well. The Planck mission (2018) gave us the exact numbers. Only 4.9\% of the universe is ordinary matter --- atoms, stars, planets, people. Dark matter accounts for 26.8\% and dark energy for 68.3\%. In practice, 95.1\% of the universe is unknown.

Imagine the universe as a chest holding 100 coins. Only 5 coins are visible. The other 95 are invisible. Of those, 27 are dark matter and 68 are dark energy. We know of dark matter that it exists (the gravitational evidence is solid), that it neither emits nor absorbs light, and that it passes through us without interacting. Roughly 100 million particles pass through your body every second. We do not know what it is made of.

A recent constraint (2025), due to Bonvin et al., has shown that dark matter falls into gravitational wells exactly as ordinary matter does. If a ``fifth force'' exists, it cannot exceed 7\% of the gravitational force.

\subsection{What is ART-12D?}
ART-12D starts from a simple idea: the universe has more dimensions than the four we perceive. These extra dimensions are compact --- so small that we do not notice them. Dark matter exists in all twelve dimensions. When it is projected into 4D, it appears as a field that fills the universe. This field is what the brain receives as consciousness.

The name comes from: \textbf{Akashic} (the universal information field), \textbf{Resonance} (because the interaction is a resonance), \textbf{Theory} (a theoretical framework) and \textbf{12D} (because I work in twelve dimensions).

The three postulates are: 1. Dark matter is a fundamental quantum field. 2. The brain is a geometric receiver. 3. The coupling is through a Yukawa interaction.

Think of a radio. It does not create the music, it picks it out of the air. The radio waves exist whether or not the radio is switched on. When the radio breaks, the music does not vanish from the universe --- only its local reception stops. In the same way, the brain works exactly like a radio. Consciousness is not generated by neurons, it is received by them. And what they receive is the field of dark matter.

\subsection{Why twelve dimensions?}
Why twelve? Why not 10, or 11, or 26? The answer comes from string theory. Imagine a guitar string. It can vibrate in several ways. Each way corresponds to a note. In the same way, a fundamental string can vibrate in several dimensions. The number of dimensions determines what kinds of vibration are possible.

The history of dimensions: Kaluza--Klein (5 dimensions, 1919--1926), superstrings (10, 1970--1980), M-theory (11, Witten, 1995), F-theory (12, Vafa, 1996), ART-12D (12, Ofiteru, 2026). Twelve is the largest number of dimensions in which string theory is mathematically consistent. More dimensions would lead to inconsistencies; fewer would lead to an incomplete theory.

\subsection{What do the compact dimensions look like?}
Imagine a garden hose. From a distance it looks one-dimensional --- it has only length. Up close it has thickness --- it is two-dimensional (the surface of the cylinder). In the same way, our space looks three-dimensional. But up close (at the Planck scale, $10^{-35}$ m) it could have eleven dimensions. The eight compact dimensions are curled up inside every point of space.

What does ``curled up'' mean? They are not places you can travel to. They are not parallel worlds you reach by travelling. They are local properties of every point. Every point of space--time has, in addition, eight small dimensions ``attached'' to it.

\subsection{How are the dimensions connected to consciousness?}
Dark matter exists in twelve dimensions. It is a scalar field $A(x,y)$, where $x$ are the 4D coordinates and $y$ are the eight compact coordinates. When we compactify the eight dimensions, the field decomposes into Kaluza--Klein modes:
\begin{equation}
A(x,y) = \sum_n A_n(x) Y_n(y)
\end{equation}
Each mode has a mass and a characteristic frequency:
\begin{equation}
m_n^2 = \frac{n^2}{R^2}
\end{equation}
The brain, through its geometric structure (the microtubules in neurons), can resonate with certain modes. When it resonates, the field is transduced into subjective experience.

Why does architecture matter rather than size? The cerebellum has 69 billion neurons (80\% of the total), yet it does not produce consciousness. It has a \textit{feedforward} architecture --- information flows in a single direction, like a wire antenna. The cortex has only 16 billion neurons, but it is the seat of consciousness. It has a \textit{recurrent} architecture --- information circulates in loops, like a well-aimed parabolic antenna. The conclusion: what counts is not the number of neurons but the geometry of the network.

\section{Methodology}
\subsection{How I worked}
I am not a professional physicist. I have no laboratory. I have no funding. What I have is a library, time, and a question that will not leave me alone. I worked in four steps. The first step was formulating the hypothesis, starting from the intuition that consciousness is received rather than generated. The second step was the formalisation, using the tools of known mathematics: Kaluza--Klein for extra dimensions, F-theory for twelve dimensions, field theory for Lagrangians, IIT for integrated information. The third step was deriving predictions, that is, what we ought to observe if the theory is correct. I selected the predictions that can be tested now. The fourth step was proposing experiments with modest resources, beginning with one costing 50,000 dollars.

\subsection{What guided me}
I held to four criteria. Internal consistency --- the equations must not contradict one another. Falsifiability --- the theory must make predictions that can be tested. Convergence --- the theory must not contradict established physics. Honesty --- I must acknowledge that it is speculative.

\subsection{What I used}
I used Kaluza--Klein (1919) for the idea of compact dimensions, F-theory (Vafa, 1996) for the twelve-dimensional framework, IIT (Tononi, 2004) for measuring integrated information, Orch OR (Penrose, 1996) for quantum collapse in microtubules, and quantum field theory for Lagrangians and equations of motion.

\subsection{My limitations}
I must be honest about the limitations. I have no experimental evidence; the theory has not been tested. I have no peer review; I have not published in scientific journals. I have no academic affiliation; I am an independent researcher. I have no funding; my resources are limited.
